\section{Métodos espectrales}
\label{sec:spectral}

Para obtener una medida gradual, representamos cada regla con un vector numérico (\emph{embedding}) derivado de su grafo de transición, de modo que la distancia entre los vectores funcione como índice de similitud.

\subsection{Eigenvalores}

Calculamos los eigenvalores de la matriz de adyacencia del grafo de transición, los ordenamos y los usamos como \emph{embedding} de la regla; dos reglas se comparan mediante la distancia coseno entre sus vectores, que mide qué tan distintos son sus espectros. Este método tiene dos limitaciones:
\begin{itemize}
    \item \textbf{Ceguera a los transitorios}: cada componente conexa está formada por árboles que llegan a un ciclo, y como los árboles solo aportan eigenvalores nulos, el espectro depende únicamente de los atractores. Las reglas 0 y 8 de ECA tienen un único atractor, al que se llega en a lo más 1 y 2 pasos, respectivamente, por lo que sus eigenvalores son idénticos.
    \item \textbf{Dependencia del dominio}: el número de eigenvalores es igual al número de nodos $N = |\mathcal{S}|^n$, por lo que solo se pueden comparar reglas con el mismo número de estados y el mismo tamaño de espacio.
\end{itemize}

En el grafo completo, a partir de $n = 8$ se obtienen entre 67 y 74 clases (Figura~\ref{fig:all-methods}). En el cociente, a partir de $n = 12$ se obtienen entre 56 y 63, y combinarlo con el grafo completo no agrega clases (77 en ambos casos, Tabla~\ref{tab:methods}). Los grupos \{0,8\}, \{1,19\}, \{2,24,46\}, \{4,12\}, \{18,126\}, \{36,72\}, \{128,136\}, \{130,152\}, \{132,140\} y \{162,168\} tienen el mismo espectro en todos los tamaños calculados: tienen los mismos ciclos y solo difieren en sus transitorios.

\subsection{Momentos espectrales}

Para resolver la dependencia del dominio proponemos usar momentos espectrales: en lugar del espectro completo, para cada paso $k = 1, \dots, K$ medimos la proporción de configuraciones que regresan a sí mismas tras $k$ pasos,
\[
  m_k = \frac{|\{x \in \mathcal{C} : F^k(x) = x\}|}{N}.
\]
La longitud del vector no depende del número de nodos sino de la ventana $K$, por lo que se pueden comparar reglas con diferente número de estados, vecinos o espacios; $K$ debe elegirse de modo que capture los ciclos relevantes de ambos grafos. Como $m_k$ solo depende de los eigenvalores, este método hereda la ceguera a los transitorios. Con $K = 60$ se obtiene un máximo de 74 clases. En el grafo completo, al aumentar $K$ se estabiliza el número de clases conforme aumenta $n$, mientras que en el cociente $K$ casi no influye, ya que los ciclos son más cortos.

\subsection{Momentos espectrales con masa de atractor}
\label{subsec:attractor_mass}

Para distinguir reglas como la 0 y la 8 es necesario medir, además de los atractores, la fuerza con la que atraen a las demás configuraciones, es decir, cuántas configuraciones llegan a cada atractor y en cuántos pasos. La idea es que cada ciclo se siga activando cada $p$ pasos, como en los momentos espectrales, pero que en cada activación se sumen también las configuraciones que ya llegaron a él. Usando la altura $h(x)$ y el periodo $p(x)$ de cada configuración, definimos
\[
  \tilde{m}_k = \frac{|\{x \in \mathcal{C} : p(x) \mid k \ \text{ y } \ h(x) \le k - p(x)\}|}{N}, \quad k = 1, \dots, K.
\]
Es decir, una configuración cuenta en el paso $k$ si su atractor se activa en ese paso y ella ya había llegado a él antes de la activación anterior: los nodos del ciclo ($h = 0$) cuentan desde la primera activación ($k = p$), los nodos con $1 \le h \le p$ desde la segunda ($k = 2p$), y así sucesivamente. Como toda configuración periódica cumple la condición cuando $p \mid k$, el vector se descompone como $\tilde{m}_k = m_k + t_k$, donde $m_k$ es el momento espectral y $t_k$ es la proporción de configuraciones transitorias que cuentan en el paso $k$. Así, el nuevo vector extiende a los momentos espectrales con la información de las cuencas y conserva su longitud $K$ independiente del dominio, con la limitación de que un ciclo con $p > K$ nunca se activa.

Por ejemplo, con 5 células ($N = 32$), la regla 0 envía todas las configuraciones a cero en un paso ($h \le 1$), mientras que en la regla 8 cada bloque de unos deja un 1 aislado que desaparece en el siguiente paso ($h \le 2$); de las 32 configuraciones, 12 llegan a cero en a lo más un paso. Sus vectores, $\tilde{m}^{(0)} = \tfrac{1}{32}(1, 32, 32, \dots)$ y $\tilde{m}^{(8)} = \tfrac{1}{32}(1, 12, 32, \dots)$, ya no son iguales, aunque sus momentos espectrales sí lo son.

Con la masa de atractor el número de clases aumenta considerablemente (Figura~\ref{fig:am-classes}): en el grafo completo, con $K = 60$ y $n \ge 6$ se obtienen entre 79 y 83 clases, y el cociente obtiene hasta 80. Las reglas 0 y 8 quedan separadas para todo $n \ge 3$, y combinando $n = 3$ y $n = 6$ se obtienen 86 clases con el grafo completo. Solo quedan juntos \{1,19\}, cuyas alturas exactas son distintas pero caen en el mismo intervalo de $p$ pasos, y \{4,12\}, que coinciden en la distribución de altura y periodo y solo difieren en la forma de los árboles. El cociente, en cambio, agrupa las reglas que difieren en componer con la traslación (\{15,51\}, \{170,204\}, \{4,12,34\}, \ldots) y solo alcanza 81 clases. Usando el grafo completo y el cociente en el mismo tamaño se obtienen 87 clases con $n = 6, 10, 12, 14, 18, 22, 24$ y $26$.

\begin{figure}[!htbp]
    \centering
    \includegraphics[width=\textwidth]{attractor_mass_classes_combined_gpu.png}
    \caption{Número de vectores de momentos espectrales con masa de atractor distintos entre las 88 reglas representantes para $n = 1, \dots, 29$ y $K = 10, 20, 30, 60$, en el grafo completo (línea continua) y en el grafo cociente (línea punteada).}
    \label{fig:am-classes}
\end{figure}

En la Figura~\ref{fig:all-methods} y en la Tabla~\ref{tab:methods} se observa la comparativa de todos los métodos. La forma canónica distingue más clases, pero los momentos espectrales con masa de atractor son el segundo mejor método y pueden producir 87 clases en total, seguidos por los eigenvalores y los momentos espectrales simples.

\begin{table}[!htbp]
\centering
\setlength{\tabcolsep}{4pt}
\begin{tabular}{lccccc}
\hline
 & \multicolumn{2}{c}{Grafo completo} & \multicolumn{2}{c}{Cociente} & Completo \\
Método & Un tamaño & Todos & Un tamaño & Todos & + cociente \\
\hline
Forma canónica & 85 & 88 & 85 & 85 & 88 \\
Eigenvalores & 74 & 77 & 63 & 68 & 77 \\
Momentos espectrales & 74 & 77 & 63 & 68 & 77 \\
Masa de atractor & 83 & 86 & 80 & 81 & 87 \\
\hline
\end{tabular}
\caption{Número máximo de clases entre las 88 reglas representantes. ``Un tamaño'': mejor $n$; ``Todos'': combina los 29 tamaños; ``Completo + cociente'': ambos grafos.}
\label{tab:methods}
\end{table}
